%% =====================================================================
%%  informe-sr.tex  --  PLANTILLA/TEST del Informe de Semana de
%%                      Representacion (tipo informe-sr)
%% ---------------------------------------------------------------------
%%  Ejercita los componentes de structures/informe-sr/ con datos de
%%  muestra: es el "test" del tipo (make plantillas). Primero la EDICION
%%  PUBLICA (revista: prosa a dos columnas, fotos grandes, sin cajas ni
%%  tablas) y luego, en paginas aparte, los componentes de la EDICION
%%  INTERNA (fichas, cajas, matriz). Los informes reales NO se escriben a
%%  mano: los genera generadores/informe-sr/main.py desde YAML
%%  y compila con scripts/build.sh.
%%
%%  Uso:  scripts/build.sh templates/informe-sr/informe-sr.tex
%% =====================================================================
\documentclass{informe-sr}
%% =====================================================================
%%  templates/informe-sr/metadata.tex  --  Datos de la PLANTILLA de prueba
%% ---------------------------------------------------------------------
%%  En los informes reales este archivo lo escribe el generador
%%  generadores/informe-sr desde datos/informe.yaml del expediente.
%% =====================================================================
\setmeta{numero}{00}
\setmeta{titulocorrido}{Informe de Semana de Representación N.º 00 · Plantilla}
\setmeta{fecha}{3 de septiembre de 2026}
\setmeta{lugar}{Ayacucho}


\begin{document}

%% ---------------------------------------------------------------- PORTADA
\portadainforme[
  kicker={Informe N.º 00 · Plantilla},
  ambito={Provincias de La Mar y Huamanga · Región Ayacucho},
  codigo={Fotografía de la Semana, 3:2},
]

%% ------------------------------------------------------- CONTENIDO + FICHA
\renewcommand{\contentsname}{Contenido}
\tableofcontents
\vspace{22pt}
\noindent\begin{minipage}[t]{0.52\linewidth}\vspace{0pt}
\isrkicker{Ficha técnica}
\isrdato{Documento}{Informe de Semana de Representación N.º 00 (plantilla)}
\isrdato{Código y versión}{ISR-00 · 2026-08 · versión 0.1 (borrador)}
\isrdato{Equipo}{\meta{autor}, \meta{cargo}\ifmetadef{distrito}{ por \meta{distrito}}{}.\\Coordinación distrital del despacho.}
\end{minipage}\hfill
\begin{minipage}[t]{0.40\linewidth}\vspace{0pt}
\isrkicker{Notas de esta versión}
{\fontsize{8.6}{11.5}\selectfont\color{colorisrtextosec}Documento de prueba del tipo documental; los textos son de muestra.\par}
\end{minipage}\par

%% ----------------------------------------------------------- PRESENTACIÓN
\clearpage
\membreteinforme
\seccion{Presentación}
\noindent\begin{minipage}[t]{0.66\linewidth}\vspace{0pt}
La Semana de Representación es el periodo en el que el diputado deja el
Palacio Legislativo para escuchar, verificar y recoger en el territorio
las demandas de la población que representa. Este informe da cuenta de
lo actuado y, sobre todo, de lo que sigue: cada pedido recibido tiene
una entidad competente, una actuación parlamentaria y un responsable de
seguimiento.
\end{minipage}\hfill
\begin{minipage}[t]{0.28\linewidth}\vspace{4pt}
\includegraphics[width=\linewidth]{\despachoRetratoBusto}\par\vspace{5pt}
{\isrtxpie Fotografía oficial del diputado.\par}
\end{minipage}\par
\vfill
\bloquefirma[1.6cm]

%% ------------------------------------------------------- RESUMEN + CRÓNICA
\clearpage
\seccion{Resumen ejecutivo}
\noindent
\isrcifra{31}{actividades en 5 días de trabajo en el territorio}{Relato cronológico}\hfill
\isrcifra{59}{pedidos y hallazgos con entidad competente}{Matriz de sistematización}\hfill
\isrcifra{25,8 h}{de actividad con horario registrado}{Relato cronológico}\hfill
\isrcifra{69}{entidades competentes identificadas}{Matriz de sistematización}\par
\vspace{16pt}
\begin{columnas}
La Semana de Representación de agosto de 2026 recorrió cuatro distritos
de la provincia de La Mar y la ciudad de Ayacucho. La agenda estuvo
dominada por la educación: la mayoría de las actividades fueron visitas
de trabajo a instituciones educativas de inicial, primaria, secundaria,
educación básica especial y alternativa, muchas de ellas en el valle de
los ríos Apurímac, Ene y Mantaro.

Los hallazgos se repiten de un colegio a otro: infraestructura precaria
o sin liquidación de obra, losas deportivas sin techar, falta de
personal administrativo y de servicio, textos escolares incompletos,
mobiliario insuficiente y conectividad que no alcanza.
\end{columnas}
\isrkicker{La ruta, día por día}
\begin{isrcrono}
  \item[24 de agosto] \textbf{Distrito de Ayna – San Francisco, provincia de La Mar.} 8 actividades; 22 pedidos sistematizados.
  \item[25 de agosto] \textbf{Distrito de Samugari, provincia de La Mar.} 10 actividades; 15 pedidos sistematizados.
\end{isrcrono}

\clearpage
\articulo{Día 1 · Lunes 24 de agosto · Ayna – San Francisco}{San Francisco:
  un hospital de seis años, aulas de triplay y la agenda del VRAEM}{Ocho
  actividades entre las 9:30 y las 20:40: la fiscalización de la obra del
  hospital, cuatro instituciones educativas, el juez de paz, los damnificados
  del huaico y la Federación de Productores de Cacao y Café.}
\begin{columnas}
La jornada comenzó con una visita inopinada a la obra del Hospital de
San Francisco, que lleva más de seis años en ejecución. Los trabajos se
reiniciaron el 23 de junio de 2026 a cargo del Consorcio Salud La Mar
II, con entrega prevista para el 19 de marzo de 2027. En la reunión con
el residente y el supervisor de obra se identificaron dos riesgos
concretos: el plan de riesgos vigente data de 2016 y no prevé los
huaicos, y las pruebas de los sistemas de agua y de los equipos
requieren 0,56 litros por segundo cuando hoy se dispone de 0,11.

El despacho oficiará a la unidad ejecutora para que remita el
cronograma valorizado, el avance físico-financiero y el plan de riesgos
actualizado; pedirá control concurrente a la Contraloría y convocará
una mesa técnica para resolver el abastecimiento de agua.
\end{columnas}
\isrcita{Una derivación sin respuesta o sin comunicación al ciudadano no
  constituye atención concluida.}{Pauta de seguimiento del despacho}{Matriz
  de sistematización, agosto de 2026}
\subseccion{Pedidos que el despacho propone atender primero}
\begin{isrpedidos}
  \item[P-011] Techado de losa; nueva infraestructura por rajaduras y paredes de triplay; acceso al agua. {\isrtxmeta UGEL La Mar; DRE Ayacucho; PRONIED.}
  \item[P-027] Falta de medicamentos, especialistas y equipos médicos básicos. {\isrtxmeta DIRESA Ayacucho; Red de Salud San Francisco.}
\end{isrpedidos}

%% ---------------------------------------------- EDICIÓN INTERNA (componentes)
\capitulo{Componentes de la edición interna}{Fichas de actividad, cajas y
  matriz: el documento de trabajo del que nace la prosa del informe público.}

\begin{fichaactividad}{id={A-01}, titulo={Visita inopinada al Hospital de San
  Francisco}, tipo={Fiscalización}, sector={salud}, hora={09:30–11:30},
  duracion={2 h}, lugar={Obra del Hospital de San Francisco},
  ubicacion={San Francisco · Ayna · La Mar},
  participantes={Residente de obra; supervisor y equipo técnico},
  resumenpedidos={2 pedidos sistematizados · P-001, P-002},
  documentos={OF-0040}, foto={slot}}
  \isrcontexto{La obra lleva más de seis años en ejecución.}
  \isrsubtitulo{Hallazgos}
  \begin{isrhallazgos}
    \item No existe un plan de riesgos actualizado: el vigente data de 2016.
    \item Para las pruebas se requieren 0,56 L/s; hoy se dispone de 0,11 L/s.
  \end{isrhallazgos}
\end{fichaactividad}

\begin{isrcaja}[nota]{Criterio de competencia}
  El diputado fiscaliza, solicita información, articula y propone
  normas; no sustituye a las entidades ejecutoras ni interfiere en
  órganos constitucionalmente autónomos.
\end{isrcaja}

\isrkicker{Matriz de sistematización (muestra)}
\begin{isrmatriz}
  P-001 & Obra con más de seis años de ejecución; sin plan de riesgos
    actualizado. & Gobierno Regional de Ayacucho; unidad ejecutora;
    Contraloría/OCI. & Oficiar a la unidad ejecutora por cronograma
    valorizado, avance y plan de riesgos. & \isrprioridad{alta} &
    \isrestado{registrado} \\
  P-027 & Falta de medicamentos, especialistas y equipos médicos básicos.
    & DIRESA Ayacucho; Red de Salud San Francisco. & Solicitar stock,
    cartera de servicios y plan de abastecimiento. & \isrprioridad{urgente}
    & \isrestado{derivado} \\
\end{isrmatriz}
\isrfuente{Matriz de sistematización de la Semana de Representación (despacho).}

\contraportadainforme
\end{document}
